\slajdid{09_1-s038}
\begin{frame}[label=09_1-s038]
\gusttekst\begin{columns}[c,onlytextwidth]\column{.23\textwidth}\centering\input{slike/tikz/s038_invertor_sa_aktivnim_opterecenjem.tex}\column{.75\textwidth}
1. slučaj kada tranzistor $T_L$ radi uvek u omskoj oblasti.\par\bigskip
Za očekivati je da ćemo dobiti slične rezultate kao i za pasivno opterećenje na izlazu.
\end{columns}\medskip
Za ulazni napon $V_I<V_{Tn,D}$ tranzistor $T_D$ (D – drive) je zakočen. Kolo je neoptrećeno $I_{Dn,D}=0=I_{Dn,L}$. Kako tranzistor $T_L$ radi u omskoj oblasti
\[I_{Dn,L}=\frac{k_{n,L}}2\bigl(2V_{DS,L}(V_{GS,L}-V_{Tn,L})-V_{DS,L}^2\bigr)=0\]
njegova radna tačka će se podesiti tako da je $V_{DS,L}=0$. To je jedina moguća radna tačka nezavisno od napona $V_{GS,L}$. Prema tome
\[V_O=V_{OH}=V_{DD}-V_{DS,L}=V_{DD}\]
\end{frame}
